\subsection {Counting Necessary Nodes (Codeforces 2074F)}
\subsection*{Problem Description}
A \textbf{quadtree} is a tree data structure in which each node has at most four children and accounts for a square-shaped region.

Formally, \textbf{for all tuples} of nonnegative integers $k, a, b \ge 0$, there exists \textbf{exactly one node} accounting for the following region:
$$[a \cdot 2^k, (a + 1) \cdot 2^k] \times [b \cdot 2^k, (b + 1) \cdot 2^k]$$

All nodes whose region is larger than $1 \times 1$ contain four children corresponding to the regions divided equally into four, and the nodes whose region is $1 \times 1$ correspond to the leaf nodes of the tree.

The pink soldiers have given you a finite region $[l_1, r_1] \times [l_2, r_2]$, where $l_i$ and $r_i$ ($l_i < r_i$) are nonnegative integers. Please find the minimum number of nodes that you must choose in order to make the union of regions accounted for by the chosen nodes \textbf{exactly} the same as the given region. Here, two sets of points are considered different if there exists a point included in one but not in the other.

\subsubsection*{Input}
Each test contains multiple test cases. The first line contains the number of test cases $t$ ($1 \le t \le 10^4$). The description of the test cases follows.

The only line of each test case contains four integers $l_1, r_1, l_2, r_2$ --- the boundaries of the region in each axis ($0 \le l_i < r_i \le 10^6$).

\subsubsection*{Output}
For each test case, output the minimum number of nodes necessary to satisfy the condition on a separate line.

\subsubsection*{Example}
\textbf{Input:}
\begin{verbatim}
5
0 1 1 2
0 2 0 2
1 3 1 3
0 2 1 5
9 98 244 353
\end{verbatim}

\textbf{Output:}
\begin{verbatim}
1
1
4
5
374
\end{verbatim}


\subsection*{Solution Approach}
Notice how the intervals on each axis correspond to those found in nodes of standard 1D segment trees. Let us find the segment tree intervals required to cover the query range on each axis individually, and call those sets of intervals $I_x$ and $I_y$.

Then, for each rectangular region formed by some $a \times b$ such that $a \in I_x$ and $b \in I_y$, the following properties hold:
\begin{enumerate}
    \item No quadtree node with side length strictly greater than $\min(|a|, |b|)$ can cover this rectangle without spilling outside the given total region.
    \item All quadtree nodes with side length no greater than $\min(|a|, |b|)$ are either completely inside the rectangle or completely outside.
\end{enumerate}

Thus, we find the following strategy to cover the region optimally:
For all rectangles $a \times b$ such that $a \in I_x$ and $b \in I_y$, cover the sub-region using the largest possible valid nodes. Since both $|a|$ and $|b|$ are powers of two (as they are segment tree nodes), the larger side length is perfectly divisible by the smaller side length. 

Therefore, we can completely tile the $a \times b$ sub-region with square nodes of side length $\min(|a|, |b|)$. The number of square nodes required to cover this specific sub-region is:
$$ \frac{\max(|a|, |b|)}{\min(|a|, |b|)} $$

Summing this value over all pairs of $a \in I_x$ and $b \in I_y$ gives the final answer. Implementing this leads to a solution with time complexity $\mathcal{O}(\log^2 X)$, where $X$ is the maximum coordinate value. While faster $\mathcal{O}(\log X)$ solutions exist, they are not strictly necessary to pass the time limits for this problem.

\subsection*{Implementation}
The following is an implementation of the logic described above. It recursively finds the optimal 1D segments for both axes and then computes the required quadtree squares by pairing them up.

\begin{lstlisting}[language=C++]
#include <bits/stdc++.h>
using namespace std;
using ll = long long;

int main() {
    ios_base::sync_with_stdio(false);
    cin.tie(NULL);
    
    int t;
    cin >> t;
    for (int i = 0; i < t; i++) {
        int l1, r1, l2, r2;
        cin >> l1 >> r1 >> l2 >> r2;
        vector<pair<ll, ll>> it1, it2;
        
        auto rec = [&](auto rec, int L, int R, int l, int r,
                       vector<pair<ll, ll>>& v) -> void {
            if (r <= L || l >= R) return;
            if (l <= L && R <= r) {
                v.emplace_back(L, R);
                return;
            }
            rec(rec, L, (L + R) / 2, l, r, v);
            rec(rec, (L + R) / 2, R, l, r, v);
        };
        
        rec(rec, 0, 1 << 25, l1, r1, it1);
        rec(rec, 0, 1 << 25, l2, r2, it2);
        
        ll ans = 0;
        
            for (auto [al, ar] : it1) {
            for (auto [bl, br] : it2) {
                ll a = ar - al;
                ll b = br - bl;
                if (a < b) swap(a, b);
                ans += a / b;
            }
        }
        cout << ans << "\n";
    }
    return 0;
}
\end{lstlisting}